\chapter{Apprentissage de représentations}\label{ch:ml:representation-learning}

Un desk dispose de plusieurs années de données de carnet d'ordres pour un instrument, mais n'en a étiqueté que quelques semaines~: celles qui ont suivi le début de la négociation
sur une nouvelle plateforme, seule période pour laquelle existe l'étiquette qui l'intéresse. Les années non étiquetées ne sont pas inutiles. Elles
montrent comment les carnets évoluent, quels états se succèdent, ce qui relève du bruit et ce qui persiste, et un modèle peut apprendre tout cela avant
d'avoir vu la moindre étiquette. Sur les séances simulées du chapitre, un encodeur pré-entraîné à prévoir les quelques
secondes suivantes du carnet lui-même, puis lu par une \omterm{def:ml:representation-learning:probe}{sonde linéaire}, surpasse tous les modèles entraînés sur le carnet brut à partir des seules étiquettes quand il y en a
300 ou 1\,000 fenêtres étiquetées~; un encodeur pré-entraîné à reconstruire son entrée n'y parvient pas. Ce chapitre construit
les objectifs auto-supervisés, mesure lesquels aident et dans quelles conditions, et explique pourquoi la reconstruction est le mauvais prétexte
pour une prévision.

\section{Représentations et autoencodeurs}

\begin{definition}[Apprentissage de représentations]\label{def:ml:representation-learning:representation}
L'\emph{apprentissage de représentations}\index{apprentissage de représentations} ajuste une application $h = e_\theta(x)$ qui transforme des entrées brutes en un
code $h$ de plus petite dimension, sur lequel les tâches ultérieures sont plus faciles~; l'application se juge à ce que ces tâches en tirent.
\end{definition}

Le chapitre~9 a défini l'\omterm{def:ml:cross-sectional-deep-models:ae}{autoencodeur}~: un encodeur, un décodeur et une perte de reconstruction. Sa leçon la plus ancienne est aussi la plus
utile ici.

\begin{proposition}[Un autoencodeur linéaire est une ACP]\label{prop:ml:represent:pca}
Parmi les encodeurs linéaires $h = Ax$ et les décodeurs $\hat x = Bh$ à code de dimension $k$, l'erreur de reconstruction
$\sum_i\lVert x_i - BAx_i\rVert^2$ est minimale lorsque $BA$ est la projection orthogonale sur le sous-espace engendré par les $k$ premières
composantes principales des données (centrées) (Baldi et Hornik, 1989).
\end{proposition}

\begin{proof}
$BAX$ est de rang au plus $k$, et d'après le théorème d'Eckart--Young, la meilleure approximation de rang $k$ de $X$ au sens de la norme de Frobenius est
sa décomposition en valeurs singulières tronquée, c'est-à-dire la projection sur les $k$ premiers vecteurs singuliers à droite, qui sont les
directions principales (livre~4, chapitres~22 et~25). Cette projection a bien la forme $BA$.
\end{proof}

Un objectif de reconstruction conserve les directions de plus forte variance, quel que soit leur rapport avec la tâche. Sur un carnet d'ordres,
la plus forte variation porte sur les tailles des niveaux profonds, qui disent peu de choses sur les secondes suivantes~; le
déséquilibre en haut du carnet, qui en dit beaucoup, ne représente qu'une petite part de la variance.

\begin{definition}[Autoencodeur débruiteur]\label{def:ml:representation-learning:dae}
Un \emph{autoencodeur débruiteur}\index{autoencodeur débruiteur} reconstruit l'entrée propre à partir d'une copie corrompue (les coefficients
masqués à zéro, un bruit gaussien ajouté), de sorte que le code doit capturer la structure de l'entrée plutôt que de la recopier (Vincent
et al., 2008).
\end{definition}

\section{Objectifs auto-supervisés}

\begin{definition}[Apprentissage auto-supervisé, pré-entraînement, ajustement fin]\label{def:ml:representation-learning:ssl}
L'\emph{apprentissage auto-supervisé}\index{apprentissage auto-supervisé} entraîne un modèle sur une tâche prétexte dont les cibles sont construites à partir des
données non étiquetées elles-mêmes~: une partie masquée de l'entrée, le futur de l'entrée, ou une autre vue du même exemple.
Le \emph{pré-entraînement}\index{pré-entraînement} est un tel ajustement, effectué avant que les étiquettes de la tâche ne soient utilisées~; \emph{ajustement fin}\index{ajustement fin}
poursuit l'entraînement du réseau pré-entraîné, avec une nouvelle tête de sortie, sur la tâche étiquetée, en général avec un petit \omterm{def:ml:trees-and-boosting:boosting}{taux
d'apprentissage}.
\end{definition}

\begin{definition}[Sonde linéaire]\label{def:ml:representation-learning:probe}
Une \emph{sonde linéaire}\index{sonde linéaire} est un modèle linéaire (ici une régression ridge) ajusté sur les codes figés
d'un encodeur pré-entraîné~; son score mesure la part de la tâche que la représentation rend accessible linéairement (Alain et
Bengio, 2017).
\end{definition}

Prévoir le futur de l'entrée elle-même est le prétexte qu'utilisent les modèles de langage, et il est disponible sur n'importe quelles données de marché~:
le carnet quelques secondes plus tard n'est pas une étiquette, c'est la suite du même flux. Le troisième encodeur du chapitre est
pré-entraîné à prédire la variation de l'instantané complet à dix niveaux deux secondes après la fin de la fenêtre. L'étiquette
qui intéresse le desk (le mid moyen sur les dix secondes suivantes) n'est jamais utilisée pendant le \omterm{def:ml:representation-learning:ssl}{pré-entraînement}.

\section{Apprentissage contrastif et augmentations pour les données de marché}

\begin{definition}[Apprentissage contrastif, augmentation de données]\label{def:ml:representation-learning:contrastive}
L'\emph{augmentation de données}\index{augmentation de données} produit de nouvelles entrées d'entraînement par des transformations qui ne doivent pas modifier
ce qui compte dans les exemples (bruit, masquage, petits décalages). L'\emph{apprentissage contrastif}\index{apprentissage contrastif} entraîne un encodeur de sorte que
deux vues augmentées d'un même exemple aient des codes similaires, et que les vues d'exemples différents aient des codes dissemblables~; avec la perte InfoNCE,
chaque vue doit retrouver son partenaire parmi les $2m - 1$ autres vues d'un lot de $m$ exemples, au moyen d'une fonction softmax appliquée aux similarités cosinus divisées par une température
(van den Oord et al., 2018~; Chen et al., 2020).
\end{definition}

Le choix des augmentations traduit l'a priori du modèle sur ce qui ne compte pas. Les modèles d'images utilisent des recadrages et des changements de couleur~;
pour un carnet d'ordres, il n'existe pas d'ensemble admis. Les augmentations du chapitre ajoutent un bruit gaussien de 0.3 écart-type et
masquent 20~\% des coefficients d'une fenêtre. Les deux sont plausibles, et toutes deux finissent par conserver les mauvaises choses~: les codes
qu'elles produisent ne sont pas meilleurs pour la prévision que ceux de l'\omterm{def:ml:cross-sectional-deep-models:ae}{autoencodeur}.

\section{Pré-entraînement, ajustement fin et sonde linéaire}

Les données du chapitre sont des séances de \texttt{firm.tape} de dix minutes~: seize non étiquetées pour le \omterm{def:ml:representation-learning:ssl}{pré-entraînement}, quatre
étiquetées, une pour l'\omterm{def:ml:trees-and-boosting:early}{arrêt précoce} et six pour les tests. Une entrée est une fenêtre de dix instantanés d'une demi-seconde portant sur dix niveaux (400
nombres)~; la cible est la variation prospective du chapitre~8 (le mid moyen sur les dix secondes suivantes par rapport au mid courant),
ici sous forme de nombre, et mesurée par son coefficient d'information sur chaque séance de test. Les trois encodeurs (un \omterm{def:ml:neural-networks-for-noisy-tabular-data:mlp}{perceptron multicouche} à 64
unités cachées et un code de dimension 16) sont pré-entraînés pendant dix \omterm{def:ml:neural-networks-for-noisy-tabular-data:backprop}{époques}~; chaque taille d'ensemble étiqueté est moyennée sur quatre
tirages, les premières fenêtres de chaque séance étiquetée.

\begin{omfigure}
\begin{tikzpicture}
\begin{semilogxaxis}[width=11cm,height=6cm,xlabel={fenêtres étiquetées},ylabel={coefficient d'information en test},
  tick label style={font=\small},label style={font=\small},xmin=220,xmax=6000,ymin=0,ymax=0.6,
  xtick={300,1000,4444},xticklabels={300,1\,000,4\,444},grid=major,grid style={gray!25},legend columns=2,
  legend cell align=left,legend style={font=\footnotesize,at={(0.5,-0.24)},anchor=north,draw=none,fill=none,
  /tikz/every even column/.append style={column sep=6pt}},table/col sep=comma]
\addplot[gray,very thick,mark=triangle*] table[x=windows,y=raw_ridge]{figdata/ml/10-representation-learning/labels.csv};
\addplot[omProp,very thick,mark=diamond*] table[x=windows,y=from_scratch]{figdata/ml/10-representation-learning/labels.csv};
\addplot[omThm!70,thick,mark=square] table[x=windows,y=autoencoder_probe]{figdata/ml/10-representation-learning/labels.csv};
\addplot[omThm,thick,dashed,mark=o] table[x=windows,y=contrastive_probe]{figdata/ml/10-representation-learning/labels.csv};
\addplot[omDef,very thick,mark=*] table[x=windows,y=predictive_probe]{figdata/ml/10-representation-learning/labels.csv};
\addplot[omDef,thick,dashed,mark=*] table[x=windows,y=predictive_fine_tuned]{figdata/ml/10-representation-learning/labels.csv};
\addplot[black!70,thick,dotted,mark=x] table[x=windows,y=features_ridge]{figdata/ml/10-representation-learning/labels.csv};
\legend{régression ridge sur la fenêtre brute,réseau entraîné depuis zéro,sonde sur autoencodeur,sonde contrastive,sonde prédictive,
prédictive ajustée finement,régression ridge sur variables construites à la main}
\end{semilogxaxis}
\end{tikzpicture}

\omcaption{Coefficient d'information en test de la prévision des dix secondes suivantes en fonction du nombre de fenêtres étiquetées (300 et 1\,000~: moyennes
sur quatre tirages~; 4\,444~: les quatre séances étiquetées). Les encodeurs pré-entraînés n'ont jamais vu d'étiquette. Données~:
\texttt{ml\_represent.compare}.}\label{fig:ml:represent:labels}
\end{omfigure}

La \Cref{fig:ml:represent:labels} montre un seul vainqueur dans le régime à peu d'étiquettes, et une seule leçon vaut partout. Avec 300 fenêtres étiquetées
(deux minutes et demie de négociation), la sonde sur l'encodeur prédictif atteint 0.22, contre 0.12 pour la régression ridge
sur la fenêtre brute et 0.18 pour le même réseau entraîné depuis zéro~; avec 1\,000 fenêtres, 0.29 contre 0.24 et 0.20. Avec
les 4\,444 fenêtres étiquetées, la régression ridge sur la fenêtre brute a rattrapé son retard (0.45 contre 0.40 pour la sonde), et seul l'\omterm{def:ml:representation-learning:ssl}{ajustement fin}
de l'encodeur prédictif se maintient à son niveau (0.46), tout en restant devant le réseau entraîné depuis zéro (0.38). Les codes de reconstruction et
contrastifs font moins bien que la fenêtre brute avec 1\,000 et 4\,444 étiquettes, et à 300 l'\omterm{def:ml:cross-sectional-deep-models:ae}{autoencodeur} fait à peine mieux
(0.15) et l'encodeur contrastif fait pire (0.05), comme la \cref{prop:ml:represent:pca} le signale~: ils conservent ce qui varie, et non ce qui
prédit.

Les variables de flux d'ordres construites à la main du chapitre~8 battent tout le reste dès qu'il y a 1\,000 étiquettes (0.50 et 0.55), et sont les
pires de toutes avec 300 (0.05)~: dix variables et une régression ridge ajustée sur deux minutes et demie d'un flux qui a rarement
bougé. Un \omterm{def:ml:representation-learning:ssl}{pré-entraînement} sans étiquette apprend à un réseau ce que la recherche en microstructure a déjà consigné~; sa valeur est
maximale là où ce savoir manque (un nouveau marché, un nouveau type de données) et où les étiquettes sont rares.

\begin{method}[Pré-entraînement sur données de marché]\label{met:ml:represent:method}
\begin{enumerate}
\item Choisir une tâche prétexte proche de la tâche visée~: prévoir le futur propre du flux l'emporte sur sa reconstruction pour une prévision.
\item Pré-entraîner sur des séances disjointes des séances étiquetées et de test~; estimer la normalisation sur les données de \omterm{def:ml:representation-learning:ssl}{pré-entraînement}.
\item Mesurer d'abord avec une \omterm{def:ml:representation-learning:probe}{sonde linéaire}~; n'ajuster finement que si la sonde est utile, à petit \omterm{def:ml:trees-and-boosting:boosting}{taux d'apprentissage} et avec \omterm{def:ml:trees-and-boosting:early}{arrêt
précoce}.
\item Rapporter la courbe d'apprentissage en données étiquetées face à la régression ridge sur l'entrée brute et à un réseau entraîné depuis zéro~: le
\omterm{def:ml:representation-learning:ssl}{pré-entraînement} n'en vaut le coût que là où les courbes se séparent.
\end{enumerate}
\end{method}

\section{Tutoriel~: apprentissage sans étiquettes}\label{tut:ml:representation-learning:tut}

\begin{tutorial}
\textbf{Objectif.} Pré-entraîner trois encodeurs sur des séances non étiquetées et mesurer leurs codes avec des sondes et un \omterm{def:ml:representation-learning:ssl}{ajustement fin} à mesure que le
nombre d'étiquettes augmente. \textbf{État final~:} \cref{fig:ml:represent:labels}.
\begin{tutsteps}
\item \textbf{L'objectif contrastif}~: deux vues augmentées, InfoNCE sur le lot.
\omcode{code/firm/represent/firm_represent.py}{71}{97}{InfoNCE et pré-entraînement contrastif.}\label{lst:ml:represent:contrastive}
\item \textbf{Prévision auto-supervisée}~: l'encodeur prédit un futur sans étiquette de son entrée.
\omcode{code/firm/represent/firm_represent.py}{100}{115}{Pré-entraînement par prévision du futur propre du flux.}\label{lst:ml:represent:predictive}
\item \textbf{Exécuter} \texttt{ml\_represent.compare(300)}, \texttt{compare(1000)}, \texttt{compare(0)} et
\texttt{fig\_represent.py}.
\end{tutsteps}
\textbf{Pour aller plus loin.} Remplacer les augmentations contrastives par un décalage temporel d'un instantané (deux fenêtres voisines prises comme
paire positive) et voir si les codes s'améliorent~; prévoir l'instantané huit secondes à l'avance au lieu de deux.
\end{tutorial}

\section{Construction~: apprentissage de représentations}\label{bld:ml:representation-learning:represent}

\begin{build}
\bfield{Rôle} Encodeurs pré-entraînés sur les données de marché non étiquetées de la firme, mesurés avant qu'un modèle ne soit construit sur eux.
\bfield{Interface} \texttt{Encoder(d\_in, code, hidden)}, \texttt{augment(x, gen, noise, mask)},
\texttt{pretrain\_dae}, \texttt{pretrain\_contrastive}, \texttt{pretrain\_predictive(X, Y, code, hidden, epochs, seed)},
\texttt{info\_nce(z1, z2, tau)}, \texttt{encode(enc, X)}, \texttt{linear\_probe(Z, y, Zt, alpha)},
\texttt{fine\_tune(enc, X, y, Xv, yv, lr, epochs, patience, seed)}, \texttt{scratch(X, y, Xv, yv, code, hidden)}.
\bfield{Règles} Données de \omterm{def:ml:representation-learning:ssl}{pré-entraînement} disjointes des données étiquetées et de test~; chaque tirage aléatoire vient de la graine de l'exécution~; un seul
thread CPU.
\bfield{Tests d'acceptation} \texttt{code/firm/represent/tests/}~: la perte InfoNCE vaut $\ln(2m-1)$ pour des codes non informatifs et
près de zéro pour des vues identiques~; le code d'un \omterm{def:ml:representation-learning:dae}{autoencodeur débruiteur} linéaire engendre le sous-espace principal dominant~; la sonde de l'encodeur prédictif
retrouve une relation injectée entre une fenêtre et son futur~; l'\omterm{def:ml:representation-learning:ssl}{ajustement fin} ne modifie pas l'encodeur
pré-entraîné passé en argument.
\bfield{Pour aller plus loin} \omterm{def:ml:representation-learning:ssl}{Pré-entraînement} par instantanés masqués avec un petit \omterm{def:ml:sequence-models-on-order-books:attention}{transformer}~; codes regroupés sur plusieurs
instruments simulés.
\end{build}

\begin{omsources}
\item G.~E.~Hinton et R.~R.~Salakhutdinov, ``Reducing the dimensionality of data with neural networks'', \emph{Science}
313, 2006.
\item P.~Vincent, H.~Larochelle, Y.~Bengio et P.-A.~Manzagol, ``Extracting and composing robust features with denoising
autoencoders'', \emph{ICML}, 2008.
\item P.~Baldi et K.~Hornik, ``Neural networks and principal component analysis: learning from examples without local
minima'', \emph{Neural Networks} 2(1), 1989.
\item A.~van den Oord, Y.~Li et O.~Vinyals, ``Representation learning with contrastive predictive coding'', arXiv, 2018;
T.~Chen, S.~Kornblith, M.~Norouzi et G.~Hinton, ``A simple framework for contrastive learning of visual representations'',
\emph{ICML}, 2020.
\item Z.~Yue et al., ``TS2Vec: towards universal representation of time series'', \emph{AAAI}, 2022.
\item G.~Alain et Y.~Bengio, ``Understanding intermediate layers using linear classifier probes'', arXiv, 2016.
\end{omsources}

\section{Exercices}

\begin{exercise}[$\star$]\label{exo:ml:representation-learning:1}
Un lot compte $m = 256$ exemples, avec deux vues chacun. Parmi combien de candidats chaque vue doit-elle retrouver son partenaire sous InfoNCE,
et quelle est la perte si les codes de l'encodeur ne portent aucune information~?
\end{exercise}

\begin{exercise}[$\star$]\label{exo:ml:representation-learning:2}
Combien de minutes de négociation représentent 300 fenêtres d'une demi-seconde chacune~? Et 4\,444~?
\end{exercise}

\begin{exercise}[$\star$]\label{exo:ml:representation-learning:3}
Un encodeur a une entrée de dimension 400, 64 unités cachées et un code de dimension 16. Combien de paramètres possède-t-il, et combien
en compte une \omterm{def:ml:representation-learning:probe}{sonde linéaire} posée sur son code~?
\end{exercise}

\begin{exercise}[$\star\star$]\label{exo:ml:representation-learning:4}
Pourquoi le code de l'\omterm{def:ml:cross-sectional-deep-models:ae}{autoencodeur} fait-il moins bien que la fenêtre brute pour la prévision, selon la \cref{prop:ml:represent:pca}~?
\end{exercise}

\begin{exercise}[$\star\star$]\label{exo:ml:representation-learning:5}
Pourquoi la régression ridge sur la fenêtre brute rattrape-t-elle la sonde prédictive lorsque les données étiquetées augmentent~?
\end{exercise}

\begin{exercise}[$\star\star$]\label{exo:ml:representation-learning:6}
\emph{Trouver la faille.} «~Nous avons pré-entraîné notre encodeur sur toutes nos données, y compris les mois de test, puisque le \omterm{def:ml:representation-learning:ssl}{pré-entraînement} n'utilise aucune
étiquette.~»
\end{exercise}

\begin{exercise}[$\star\star\star$]\label{exo:ml:representation-learning:7}
\emph{Programmation.} Pré-entraîner l'encodeur prédictif pour prévoir la variation de l'instantané huit secondes à l'avance (\texttt{AHEAD = 16})
et relancer \texttt{compare(1000)}. La sonde s'améliore-t-elle~?
\end{exercise}

\begin{exercise}[$\star\star\star$]\label{exo:ml:representation-learning:8}
Montrer que la perte InfoNCE d'un lot est l'entropie croisée d'une classification à $(2m - 1)$ classes et qu'elle vaut
$\ln(2m - 1)$ lorsque toutes les similarités sont égales.
\end{exercise}

\section{Problème~: apprentissage sans étiquettes}

\begin{problem}[{Problème du week-end --- trois prétextes et quelques étiquettes}]\label{pb:ml:representation-learning:1}
Les sessions du chapitre~: seize sans étiquette, quatre étiquetées, une pour l'\omterm{def:ml:trees-and-boosting:early}{arrêt précoce} (early stopping), six pour les tests.

\medskip\noindent\textbf{Partie I --- Les prétextes.}
\begin{enumerate}
\item Qu'apprend chacun des trois encodeurs à faire~?
\item Lequel d'entre eux est le plus proche de la tâche, et pourquoi~?
\item Que prédit la \cref{prop:ml:represent:pca} pour l'\omterm{def:ml:cross-sectional-deep-models:ae}{autoencodeur}~?
\item Quel a priori les augmentations contrastives encodent-elles~?
\end{enumerate}

\medskip\noindent\textbf{Partie II --- Peu d'étiquettes.}
\begin{enumerate}[resume]
\item Quels scores obtiennent la sonde prédictive, la régression ridge sur la fenêtre brute et le réseau entraîné sans \omterm{def:ml:representation-learning:ssl}{pré-entraînement} avec 300 fenêtres étiquetées~?
\item Et avec 1\,000~?
\item Comment se comparent les sondes issues de l'\omterm{def:ml:cross-sectional-deep-models:ae}{autoencodeur} et de l'\omterm{def:ml:representation-learning:contrastive}{apprentissage contrastif}~?
\item Pourquoi les variables construites à la main sont-elles les moins bonnes à 300 et les meilleures à 1\,000~?
\end{enumerate}

\medskip\noindent\textbf{Partie III --- Beaucoup d'étiquettes.}
\begin{enumerate}[resume]
\item Quels scores obtiennent les modèles avec les 4\,444 fenêtres étiquetées~?
\item Pourquoi l'\omterm{def:ml:representation-learning:ssl}{ajustement fin} (fine-tuning) aide-t-il ici alors que la sonde n'aide pas~?
\item Quel est le gain du modèle ajusté par rapport au réseau entraîné sans \omterm{def:ml:representation-learning:ssl}{pré-entraînement}~?
\item Que peut-on attendre avec quarante sessions étiquetées~?
\end{enumerate}

\medskip\noindent\textbf{Partie IV --- Le verdict.}
\begin{enumerate}[resume]
\item Énoncer le \emph{résultat nommé}~: la taille de l'échantillon étiqueté en deçà de laquelle le \omterm{def:ml:representation-learning:ssl}{pré-entraînement} aide, et l'IC de la sonde comparé à celui
du modèle entraîné sans \omterm{def:ml:representation-learning:ssl}{pré-entraînement} à cette taille.
\item Quand le \omterm{def:ml:representation-learning:ssl}{pré-entraînement} vaut-il son coût sur un desk~?
\item Comment choisir les augmentations pour les carnets d'ordres~?
\item Que faut-il tenir à l'écart des données de \omterm{def:ml:representation-learning:ssl}{pré-entraînement}~?
\item Comment pré-entraîner sur plusieurs instruments~?
\item Que révèle une \omterm{def:ml:representation-learning:probe}{sonde linéaire} (linear probe) que l'\omterm{def:ml:representation-learning:ssl}{ajustement fin} ne révèle pas~?
\item Où se situe un modèle de langage pré-entraîné dans ce cadre (chapitre~14)~?
\item En une phrase~: qu'est-ce qui rend un prétexte utile~?
\end{enumerate}
\end{problem}

\section{Questions d'entretien}

\begin{interviewq}[$\star$ \iqroles{ingénieur ML}]\label{iq:ml:representation-learning:1}
Quelle est la différence entre une \omterm{def:ml:representation-learning:probe}{sonde linéaire} et l'\omterm{def:ml:representation-learning:ssl}{ajustement fin}, et quand utilisez-vous chacun~?
\end{interviewq}

\begin{interviewq}[$\star\star$ \iqroles{ingénieur ML, chercheur}]\label{iq:ml:representation-learning:2}
Expliquez l'\omterm{def:ml:representation-learning:contrastive}{apprentissage contrastif}. Quelles augmentations utiliseriez-vous pour des données de prix intrajournalières~?
\end{interviewq}

\begin{interviewq}[$\star\star$ \iqroles{chercheur}]\label{iq:ml:representation-learning:3}
Vous disposez de dix ans de données de tick non étiquetées et de trois mois d'étiquettes pour une nouvelle stratégie. Comment utilisez-vous les dix ans~?
\end{interviewq}

\begin{interviewq}[$\star\star$ \iqroles{chercheur, ingénieur ML}]\label{iq:ml:representation-learning:4}
Pourquoi le code d'un \omterm{def:ml:cross-sectional-deep-models:ae}{autoencodeur} peut-il être inutile pour la prédiction alors qu'il reconstruit bien l'entrée~?
\end{interviewq}

\begin{interviewq}[$\star\star$ \iqroles{ingénieur ML}]\label{iq:ml:representation-learning:5}
Comment évaluez-vous une représentation pré-entraînée sans vous engager dans un modèle aval~?
\end{interviewq}

\begin{interviewq}[$\star\star\star$ \iqroles{chercheur}]\label{iq:ml:representation-learning:6}
Montrez que le meilleur \omterm{def:ml:cross-sectional-deep-models:ae}{autoencodeur} linéaire à code de dimension $k$ projette sur les $k$ premières composantes principales.
\end{interviewq}
